\documentclass[10pt,a4paper]{amsart}
\usepackage{amsmath,amssymb,amsthm,graphicx}
% Retained preamble for the Stacks Project chapter morphisms.tex (Morphisms of Schemes).
% The chapter file pinned-original-source.tex carries only \input{preamble} on its first line, so no
% preamble text lives inside the chapter itself.  The Stacks Project uses one single shared preamble for
% every chapter; the snapshot stages it verbatim in sources/stacks-chow/preamble.tex.  The block below is
% copied mechanically, line for line, from that file's lines 14-18 (the Xy-pic setup) and lines 154-217
% (the native theorem declarations together with the mathematical macros).  Nothing was added, removed or
% redefined.  The wrapper class and the amsmath/amssymb/amsthm/graphicx base packages are supplied by the
% assembler, so the class/\externaldocument part of the shared preamble is not repeated here.
% For commutative diagrams we use Xy-pic
\usepackage[all]{xy}

% We use 2cell for 2-commutative diagrams.
\xyoption{2cell}
% Theorem environments.
%
\theoremstyle{plain}
\newtheorem{theorem}[subsection]{Theorem}
\newtheorem{proposition}[subsection]{Proposition}
\newtheorem{lemma}[subsection]{Lemma}

\theoremstyle{definition}
\newtheorem{definition}[subsection]{Definition}
\newtheorem{example}[subsection]{Example}
\newtheorem{exercise}[subsection]{Exercise}
\newtheorem{situation}[subsection]{Situation}

\theoremstyle{remark}
\newtheorem{remark}[subsection]{Remark}
\newtheorem{remarks}[subsection]{Remarks}

\numberwithin{equation}{subsection}

% Macros
%
\def\lim{\mathop{\mathrm{lim}}\nolimits}
\def\colim{\mathop{\mathrm{colim}}\nolimits}
\def\Spec{\mathop{\mathrm{Spec}}}
\def\Hom{\mathop{\mathrm{Hom}}\nolimits}
\def\Ext{\mathop{\mathrm{Ext}}\nolimits}
\def\SheafHom{\mathop{\mathcal{H}\!\mathit{om}}\nolimits}
\def\SheafExt{\mathop{\mathcal{E}\!\mathit{xt}}\nolimits}
\def\Sch{\mathit{Sch}}
\def\Mor{\mathop{\mathrm{Mor}}\nolimits}
\def\Ob{\mathop{\mathrm{Ob}}\nolimits}
\def\Sh{\mathop{\mathit{Sh}}\nolimits}
\def\NL{\mathop{N\!L}\nolimits}
\def\CH{\mathop{\mathrm{CH}}\nolimits}
\def\proetale{{pro\text{-}\acute{e}tale}}
\def\etale{{\acute{e}tale}}
\def\QCoh{\mathit{QCoh}}
\def\Ker{\mathop{\mathrm{Ker}}}
\def\Im{\mathop{\mathrm{Im}}}
\def\Coker{\mathop{\mathrm{Coker}}}
\def\Coim{\mathop{\mathrm{Coim}}}

% Boxtimes
%
\DeclareMathSymbol{\boxtimes}{\mathbin}{AMSa}{"02}

%
% Macros for moduli stacks/spaces
%
\def\QCohstack{\mathcal{QC}\!\mathit{oh}}
\def\Cohstack{\mathcal{C}\!\mathit{oh}}
\def\Spacesstack{\mathcal{S}\!\mathit{paces}}
\def\Quotfunctor{\mathrm{Quot}}
\def\Hilbfunctor{\mathrm{Hilb}}
\def\Curvesstack{\mathcal{C}\!\mathit{urves}}
\def\Polarizedstack{\mathcal{P}\!\mathit{olarized}}
\def\Complexesstack{\mathcal{C}\!\mathit{omplexes}}
% \Pic is the operator that assigns to X its picard group, usage \Pic(X)
% \Picardstack_{X/B} denotes the Picard stack of X over B
% \Picardfunctor_{X/B} denotes the Picard functor of X over B
\def\Pic{\mathop{\mathrm{Pic}}\nolimits}
\def\Picardstack{\mathcal{P}\!\mathit{ic}}
\def\Picardfunctor{\mathrm{Pic}}
\def\Deformationcategory{\mathcal{D}\!\mathit{ef}}

% --- appended 2026-10-12 (Chief): the project own definitions of the reference/slogan/history
% display environments, copied mechanically from the shared Stacks preamble lines 8-12
% (\usepackage{verbatim} plus the three \newenvironment lines, which the chapter preamble keeps
% verbatim from the project preamble). Nothing added, removed or redefined: these environments
% carry no mathematical content and the project itself typesets them as comments. Without these
% three lines a counted statement whose own text uses them fails with Environment slogan undefined.
\usepackage{verbatim}
\newenvironment{reference}{\comment}{\endcomment}
\newenvironment{slogan}{\comment}{\endcomment}
\newenvironment{history}{\comment}{\endcomment}
\usepackage{cleveref}
\crefname{originallemma}{original lemma}{original lemmas}
\Crefname{originallemma}{Original lemma}{Original lemmas}
\makeatletter
\newlabel{algebra-lemma-integral-overring-surjective}{{00GQ}{source}{}{sourcepointer-1}{}}
\newlabel{algebra-lemma-integral-overring-surjective@cref}{{[originallemma][0][]00GQ}{[1][0][]source}{}{}{}}
\makeatother
\crefname{originallemma}{original lemma}{original lemmas}
\Crefname{originallemma}{Original lemma}{Original lemmas}
\makeatletter
\newlabel{lemma-colimit-inherits}{{0CN8}{source}{}{sourcepointer-2}{}}
\newlabel{lemma-colimit-inherits@cref}{{[originallemma][0][]0CN8}{[1][0][]source}{}{}{}}
\makeatother
\crefname{originallemma}{original lemma}{original lemmas}
\Crefname{originallemma}{Original lemma}{Original lemmas}
\makeatletter
\newlabel{lemma-integral-universally-closed}{{01WM}{source}{}{sourcepointer-3}{}}
\newlabel{lemma-integral-universally-closed@cref}{{[originallemma][0][]01WM}{[1][0][]source}{}{}{}}
\makeatother
\crefname{originallemma}{original lemma}{original lemmas}
\Crefname{originallemma}{Original lemma}{Original lemmas}
\makeatletter
\newlabel{lemma-square-and-cube}{{0CNB}{source}{}{sourcepointer-4}{}}
\newlabel{lemma-square-and-cube@cref}{{[originallemma][0][]0CNB}{[1][0][]source}{}{}{}}
\makeatother
\crefname{originallemma}{original lemma}{original lemmas}
\Crefname{originallemma}{Original lemma}{Original lemmas}
\makeatletter
\newlabel{lemma-universal-homeomorphism}{{04DF}{source}{}{sourcepointer-5}{}}
\newlabel{lemma-universal-homeomorphism@cref}{{[originallemma][0][]04DF}{[1][0][]source}{}{}{}}
\makeatother
\crefname{originallemma}{original lemma}{original lemmas}
\Crefname{originallemma}{Original lemma}{Original lemmas}
\makeatletter
\newlabel{lemma-universally-injective}{{01S4}{source}{}{sourcepointer-6}{}}
\newlabel{lemma-universally-injective@cref}{{[originallemma][0][]01S4}{[1][0][]source}{}{}{}}
\makeatother
\title{A ring extension gives a universal homeomorphism with isomorphisms on residue fields exactly when every finite set sits in a tower with square and cube in the previous stage}
\author{The Stacks Project Authors}
\date{stacks-project:morphisms.tex:local-snapshot-20261012; GNU FDL 1.2 or later}
\begin{document}
\maketitle
\noindent\textit{Editorial note.} All delivered mathematical text is the author's own, copied from the pinned Stacks Project chapter file pinned-original-source.tex; no statement or proof step was rewritten, reordered or removed.\par\medskip
\noindent\textit{Editorial note.} Following the multi-claim convention for this source, source\_url is the source-level URL https://stacks.math.columbia.edu/ and the source label proposition-universal-homeomorphism-equal-residue-fields is carried in claim\_key and here rather than in a per-claim tag URL.\par\medskip
\noindent\textit{Editorial note.} The preamble is the chapter's shared Stacks Project preamble, retained as sources/stacks-morphisms/retained-preamble.tex: a line-for-line excerpt of the shared Stacks preamble staged with this snapshot. No macro was invented or redefined, and this card adds no preamble text of its own.\par\medskip
\noindent\textit{Editorial note.} Statement span 11712-11724 and proof span 11726-11782 come from the halfA pool card operations/pool-cards/stacks-morphisms-halfA.json and were re-read in the pinned chapter file before this card was written.\par\medskip
\noindent\textit{Editorial note.} The proof uses the author's transfinite indexing of the chain of subalgebras and the author's phrase about the number of ordinals needed; both are preserved, and no index was normalised.\par\medskip
\noindent\textit{Editorial note.} The displayed description of the base changed algebra as a quotient of a ring cut out by two relations is the author's own; one relation between the square and the cube in the base is used as written.\par\medskip
\noindent\textit{Editorial note.} 6 cross-references occur inside the counted spans; they are listed in external\_references as navigation pointers only.\par\medskip
\noindent\textit{Allowed standard prerequisites (editorial).}\begin{itemize}
\item The square and cube lemma for a ring extension: an element outside the subring with two consecutive powers inside can be replaced by one whose square and cube lie inside. Cited at proof lines 11730-11733, this being the lane's card L34015 (lemma-nth-and-nplusone-implies-square-and-cube). Supplies the element adjoined at each successor stage of the transfinite construction.
\item Behaviour of the properties at hand under filtered colimits of rings: a colimit inherits the property, and both universal closedness and inducing isomorphisms on residue fields are stable under composition. Cited at proof lines 11753-11757. Reduces the implication from the tower condition to the single-generator case.
\item An integral ring extension gives a universally closed and surjective morphism of spectra, cited from the same chapter and the algebra chapter. Cited at proof lines 11761-11764. Provides universal closedness and surjectivity in the single-generator case.
\item The characterizations of a universal homeomorphism and of universal injectivity, cited from the same chapter. Cited at proof lines 11776-11778. Converts bijectivity on points plus isomorphisms on residue fields plus universal closedness into the assertion of (1).
\end{itemize}\medskip
\ifdefined\hypertarget\hypertarget{sourcepointer-1}{}\fi
\noindent\textit{Source reference.} Original lemma 00GQ: focused-seeds/AGP/sources/stacks-snapshot-20261011/tags.txt, source lines 616--616. Same-chapter reference used in the retained author text at source line(s) 11521, 11770. It is supplied only as a navigation pointer so that the retained text compiles; the counted item delivers the author's own statement and the author's complete proof, and no part of the referenced item is copied in or substituted for the counted argument.\par\medskip
\ifdefined\hypertarget\hypertarget{sourcepointer-2}{}\fi
\noindent\textit{Source reference.} Original lemma 0CN8: focused-seeds/AGP/sources/stacks-snapshot-20261011/tags.txt, source lines 15883--15883. Same-chapter reference used in the retained author text at source line(s) 11760, 12173. It is supplied only as a navigation pointer so that the retained text compiles; the counted item delivers the author's own statement and the author's complete proof, and no part of the referenced item is copied in or substituted for the counted argument.\par\medskip
\ifdefined\hypertarget\hypertarget{sourcepointer-3}{}\fi
\noindent\textit{Source reference.} Original lemma 01WM: focused-seeds/AGP/sources/stacks-snapshot-20261011/tags.txt, source lines 2487--2487. Same-chapter reference used in the retained author text at source line(s) 10976, 11172, 11375, 11380, 11512, 11552, 11624, 11672, 11769, 14152, 14165, 14544. It is supplied only as a navigation pointer so that the retained text compiles; the counted item delivers the author's own statement and the author's complete proof, and no part of the referenced item is copied in or substituted for the counted argument.\par\medskip
\ifdefined\hypertarget\hypertarget{sourcepointer-4}{}\fi
\noindent\textit{Source reference.} Original lemma 0CNB: focused-seeds/AGP/sources/stacks-snapshot-20261011/tags.txt, source lines 15886--15886. Same-chapter reference used in the retained author text at source line(s) 11667, 11732, 11808. It is supplied only as a navigation pointer so that the retained text compiles; the counted item delivers the author's own statement and the author's complete proof, and no part of the referenced item is copied in or substituted for the counted argument.\par\medskip
\ifdefined\hypertarget\hypertarget{sourcepointer-5}{}\fi
\noindent\textit{Source reference.} Original lemma 04DF: focused-seeds/AGP/sources/stacks-snapshot-20261011/tags.txt, source lines 5725--5725. Same-chapter reference used in the retained author text at source line(s) 11301, 11487, 11529, 11568, 11780, 11886, 14718, 14908. It is supplied only as a navigation pointer so that the retained text compiles; the counted item delivers the author's own statement and the author's complete proof, and no part of the referenced item is copied in or substituted for the counted argument.\par\medskip
\ifdefined\hypertarget\hypertarget{sourcepointer-6}{}\fi
\noindent\textit{Source reference.} Original lemma 01S4: focused-seeds/AGP/sources/stacks-snapshot-20261011/tags.txt, source lines 2325--2325. Same-chapter reference used in the retained author text at source line(s) 1676, 11525, 11564, 11702, 11781, 11892, 14730. It is supplied only as a navigation pointer so that the retained text compiles; the counted item delivers the author's own statement and the author's complete proof, and no part of the referenced item is copied in or substituted for the counted argument.\par\medskip
\begin{proposition}
\label{proposition-universal-homeomorphism-equal-residue-fields}
Let $A \subset B$ be a ring extension. The following are equivalent
\begin{enumerate}
\item $\Spec(B) \to \Spec(A)$ is a universal homeomorphism inducing
isomorphisms on residue fields, and
\item every finite subset $E \subset B$ is contained in an extension
$$
A[b_1, \ldots, b_n] \subset B
$$
such that $b_i^2, b_i^3 \in A[b_1, \ldots, b_{i - 1}]$ for $i = 1, \ldots, n$.
\end{enumerate}
\end{proposition}

\begin{proof}
Assume (1). Using transfinite recursion we construct for each ordinal $\alpha$
an $A$-subalgebra $B_\alpha \subset B$ as follows. Set $B_0 = A$. If $\alpha$
is a limit ordinal, then we set $B_\alpha = \colim_{\beta < \alpha} B_\beta$.
If $\alpha = \beta + 1$, then either
$B_\beta = B$ in which case we set $B_\alpha = B_\beta$ or
$B_\beta \not = B$, in which case we apply Lemma \ref{lemma-square-and-cube}
to choose a $b_\alpha \in B$, $b_\alpha \not \in B_\beta$ with
$b_\alpha^2, b_\alpha^3 \in B_\beta$
and we set $B_\alpha = B_\beta[b_\alpha] \subset B$.
Clearly, $B = \colim B_\alpha$ (in fact $B = B_\alpha$ for some
ordinal $\alpha$ as one sees by looking at cardinalities).
We will prove, by transfinite induction, that (2) holds for
$A \to B_\alpha$ for every ordinal $\alpha$. It is clear for
$\alpha = 0$. Assume the statement holds for every $\beta < \alpha$
and let $E \subset B_\alpha$ be a finite subset.
If $\alpha$ is a limit ordinal, then
$B_\alpha = \bigcup_{\beta < \alpha} B_\beta$ and we see
that $E \subset B_\beta$ for some $\beta < \alpha$ which
proves the result in this case. If $\alpha = \beta + 1$,
then $B_\alpha = B_\beta[b_\alpha]$. Thus any $e \in E$
can be written as a polynomial $e = \sum d_{e, i}b_\alpha^i$
with $d_{e, i} \in B_\beta$. Let $D \subset B_\beta$
be the set $D = \{d_{e, i}\} \cup \{b_\alpha^2, b_\alpha^3\}$.
By induction assumption
there exists an $A$-subalgebra $A[b_1, \ldots, b_n] \subset B_\beta$
as in the statement of the lemma containing $D$.
Then $A[b_1, \ldots, b_n, b_\alpha] \subset B_\alpha$
is an $A$-subalgebra of $B_\alpha$ as in the statement
of the lemma containing $E$.

\medskip\noindent
Assume (2). Write $B = \colim B_\lambda$ as the colimit of its
finite $A$-subalgebras.
By Lemma \ref{lemma-colimit-inherits} it suffices to show
that $\Spec(B_\lambda) \to \Spec(A)$ is a
universal homeomorphism inducing isomorphisms on residue fields.
Compositions of universally closed morphisms are universally closed
and the same thing for morphisms which induce isomorphisms on residue fields.
Thus it suffices to show that if $A \subset B$ and $B$ is generated
by a single element $b$ with $b^2, b^3 \in A$, then (1) holds.
Such an extension is integral and hence $\Spec(B) \to \Spec(A)$ is
universally closed and surjective
(Lemma \ref{lemma-integral-universally-closed} and
Algebra, Lemma \ref{algebra-lemma-integral-overring-surjective}).
Note that $(b^2)^3 = (b^3)^2$ in $A$.
For any ring map $\varphi : A \to K$ to a field $K$ we see that there
exists a $\lambda \in K$ with $\varphi(b^2) = \lambda^2$ and
$\varphi(b^3) = \lambda^3$. Namely, $\lambda = 0$ if $\varphi(b^2) = 0$
and $\lambda = \varphi(b^3)/\varphi(b^2)$ if not. Thus
$B \otimes_A K$ is a quotient of $K[x]/(x^2 - \lambda^2, x^3 - \lambda^3)$.
This ring has exactly one prime with residue field $K$.
This implies that $\Spec(B) \to \Spec(A)$ is bijective and induces
isomorphisms on residue fields. Combined with universal closedness
this shows (1) is true, see Lemmas \ref{lemma-universal-homeomorphism}
and \ref{lemma-universally-injective}.
\end{proof}

\end{document}
